\documentclass[11pt]{article}
\usepackage[a4paper,margin=1in]{geometry}
\usepackage{amsmath,amssymb,mathtools,bm}
\usepackage{enumitem,booktabs,array}
\usepackage{tcolorbox}
\usepackage{hyperref}
\usepackage[protrusion=true,expansion=false]{microtype}
\hypersetup{colorlinks=true,linkcolor=blue,urlcolor=blue,citecolor=blue}
\setlist[itemize]{topsep=2pt,itemsep=2pt}
\setlist[enumerate]{topsep=2pt,itemsep=2pt}
\newtcolorbox{keybox}{colback=blue!3!white,colframe=blue!40!black,boxrule=0.4pt,arc=1mm}
\newtcolorbox{alertbox}{colback=red!3!white,colframe=red!40!black,boxrule=0.4pt,arc=1mm}
\newtcolorbox{answerbox}{colback=green!3!white,colframe=green!40!black,boxrule=0.4pt,arc=1mm}
\newtcolorbox{drillbox}{colback=yellow!5!white,colframe=orange!60!black,boxrule=0.4pt,arc=1mm}
\newcommand{\EV}{\mathbb{E}}
\newcommand{\Var}{\mathrm{Var}}
\newcommand{\Cov}{\mathrm{Cov}}
\newcommand{\blank}[1]{\underline{\hspace{#1}}}

\title{W11-02: Hong \& Kacperczyk (2009, JFE) \\
\large ``The Price of Sin: The Effects of Social Norms on Markets'' --- Active Recall Workbook}
\author{Banking \& Risk Management --- Exam Prep}
\date{\today}
\begin{document}\maketitle

\begin{keybox}
\textbf{Paper in one sentence.} ``Sin'' stocks --- publicly traded firms in alcohol, tobacco, and gaming --- earn risk-adjusted abnormal returns of $\sim 2.5$--$3$\% per year 1965--2006 and are less held by norm-constrained investors (mutual funds, pensions) and less covered by analysts --- a pricing premium that compensates for social-norm-based exclusion.

\textbf{Placement.} The canonical ``neglected stocks'' result in the ESG-pricing literature. Complements Chava (2014) and Sharfman-Fernando (2008) on environmental performance and cost of capital, from the \emph{exclusion} perspective rather than the \emph{performance} one.
\end{keybox}

\section*{Round 1: Complete Walk-Through}

\subsection*{1.1 Sin classification}
Hong-Kacperczyk (HK) define sin stocks as firms in alcohol (SIC), tobacco, and gaming (casinos). These industries face explicit exclusions by many institutional investors (pension funds, religious funds, some mutual funds).

\subsection*{1.2 Return tests}
Fama-French-Carhart 4-factor alphas on long portfolios of sin stocks and on long-sin / short-control portfolios. 1965--2006 monthly.

\subsection*{1.3 Holdings and coverage}
\begin{itemize}
\item Mutual fund and pension holdings of sin stocks: significantly lower than comparable non-sin firms.
\item Hedge funds and individual investors hold them more than proportional.
\item Analyst coverage of sin stocks: lower.
\item Implication: sin stocks are \emph{neglected} by norm-constrained institutions.
\end{itemize}

\subsection*{1.4 Headline results}
\begin{itemize}
\item Sin stocks earn $\sim 2.5$--$3$ pp annualised 4-factor alpha.
\item Alpha is robust to additional controls (litigation risk, industry factors).
\item Neglect and return premium are cross-sectionally correlated.
\end{itemize}

\subsection*{1.5 Interpretation}
\begin{itemize}
\item Norm-based exclusion segments the market: sin stocks held by a smaller set of risk-bearers.
\item Under Merton (1987) segmented-market logic, limited risk-sharing generates an equilibrium premium.
\item Quantitative mapping to social-norm intensity works well.
\end{itemize}

\subsection*{1.6 Caveats}
\begin{alertbox}
\begin{itemize}
\item Sin classification is discrete and may miss related firms.
\item Litigation risk (tobacco) could alone account for part of the alpha.
\item Out-of-sample: post-2006 performance has been mixed, suggesting the premium may decay as norms shift.
\end{itemize}
\end{alertbox}

\section*{Round 2: Level 1 Fill-in}
\begin{enumerate}
\item Sin industries: \blank{1cm}, \blank{1cm}, and \blank{1cm}.
\item Alpha magnitude: $\sim$\blank{1cm}--\blank{1cm} pp/yr.
\item Factor model: Fama-French-\blank{1cm} 4-factor.
\item Institutional holdings: \blank{1cm} (higher/lower) for sin stocks.
\item Theoretical mapping: \blank{1cm} (1987) segmented-market model.
\end{enumerate}

\section*{Round 3: Level 2 Prompts}
\begin{enumerate}
\item Define sin stocks and describe the excluded-investor population.
\item Derive the intuition for a Merton-style segmentation premium.
\item Why could litigation-risk explanation be partial?
\item Discuss post-2006 evidence on the sin premium.
\end{enumerate}

\section*{Round 4: Minimal Scaffolding}
From a blank page: (i) sin definition, (ii) 4-factor test, (iii) holdings/coverage, (iv) magnitude, (v) caveats.

\section*{Round 5: Blank Page}
Essay: social norms and equilibrium asset pricing --- the sin premium and its limits.

\vspace{6pt}
\begin{drillbox}
\textbf{Speed Drill.} (1) Three sin industries. (2) Alpha magnitude. (3) Factor model. (4) Institutional holding direction. (5) Theoretical anchor.
\end{drillbox}

\section*{Quick Reference \& Answer Key}
\begin{answerbox}
\begin{itemize}
\item \textbf{Sin.} Alcohol, tobacco, gaming.
\item \textbf{Alpha.} $\sim$2.5--3 pp/yr (4-factor).
\item \textbf{Holdings.} Lower by norm-constrained institutions.
\item \textbf{Theory.} Merton (1987) segmentation.
\end{itemize}
\end{answerbox}

\begin{keybox}
\textbf{30-second exam pitch.} Hong \& Kacperczyk (2009) provide the classic evidence that social norms affect asset prices. Public firms in alcohol, tobacco, and gaming --- ``sin stocks'' --- are excluded by many pension funds, religious funds, and some mutual funds. HK show that such stocks earn $\sim$2.5--3 pp annualised 4-factor alpha 1965--2006 and are also less covered by analysts and more concentrated in hedge-fund and individual-investor portfolios. Under Merton's (1987) segmented-market logic, a subset of investors bearing the residual risk demand extra return --- a norm-driven premium. The paper established the norm-pricing channel in the ESG literature and provided the template for later work on climate-and-sin exclusion effects, though post-2006 evidence suggests the premium may evolve as norms and institutional exclusion patterns shift.
\end{keybox}

\end{document}
